\section{Conclusion}
\label{sec:conclusion}

We presented \fullname (\ours), the first continual learning framework for knowledge graphs that is both natively multimodal and modality-aware.
\ours integrates structural (R-GCN), textual (BiomedBERT), and molecular (Morgan fingerprint MLP) encoders through cross-modal attention, and protects previously learned knowledge via modality-aware EWC with per-modality Fisher matrices and multimodal memory replay with K-means diverse selection.
On the PrimeKG-CL benchmark, \ours achieves AP\,=\,$0.063$ on link prediction ($3.5\times$ Joint Training, $1.6\times$ LKGE) and AP\,=\,$0.431$ on node classification ($1.16\times$ Joint Training), establishing new state-of-the-art performance for continual biomedical KG learning.

\paragraph{Limitations.}
Our work has several limitations that warrant transparency.
First, ablation studies (modality ablation, decoder ablation, lambda sensitivity) have only been conducted at smoke-test scale; full-scale ablations are needed to confirm component contributions.
Second, the decoder confound between DistMult (\ours) and TransE (baselines) means that some of the observed AP advantage may stem from the scoring function rather than the continual learning mechanisms.
Third, our benchmark uses only two temporal snapshots ($t_0$, $t_1$); evaluating over longer time horizons with more snapshots would better characterize long-term forgetting behavior.
Finally, the candidate filtering in DistMult evaluation (restricting to type-compatible entities) is not directly comparable to the all-entity evaluation used by TransE baselines, inflating certain task-specific metrics.

\paragraph{Future Work.}
We identify three directions for future investigation:
(1)~full-scale ablation experiments with controlled decoder comparisons to isolate the contributions of multimodal encoding, MA-EWC, and diverse replay;
(2)~extension to multi-hop reasoning evaluation, where multimodal features may provide even stronger advantages for complex inference chains; and
(3)~integration of additional databases (DrugBank, UMLS) and a third temporal snapshot ($t_2$, February 2026) to stress-test continual learning over longer horizons.
